\begin{table}[t]
    \centering
    \small
    \begin{minipage}[t]{0.47\textwidth}
    \centering
    \textbf{(a) Single-night reliability and disattenuation}\\[2pt]
    \resizebox{\ifdim\width>\linewidth\linewidth\else\width\fi}{!}{%
\begin{tabular}{@{}lccc@{}}
        \toprule
        \textbf{Metric} & \textbf{ICC(2,1)} & \textbf{obs.\ $r$} & \textbf{true $r$} \\
        \midrule
        \stage & $\mathbf{0.775}$ & $+0.026$ & $\mathbf{+0.029}$ \\
        \se    & $0.562$ & $+0.090$ & $+0.109$ \\
        \waso  & $0.517$ & $+0.109$ & $+0.144$ \\
        \tst   & $\mathbf{0.225}$ & $+0.163$ & $\mathbf{+0.228}$ \\
        \bottomrule
    \end{tabular}%
}
    \end{minipage}
    \hfill
    \begin{minipage}[t]{0.44\textwidth}
    \centering
    \textbf{(b) The measurement-only prediction}\\[2pt]
    \resizebox{\ifdim\width>\linewidth\linewidth\else\width\fi}{!}{%
\begin{tabular}{@{}lccc@{}}
        \toprule
        \textbf{Ratio} & \textbf{predicted} & \textbf{observed} & \textbf{obs./pred.} \\
        \midrule
        \stage $\div$ \tst & $\mathbf{1.86\times}$ & $\mathbf{0.16\times}$ & $0.09$ \\
        \stage $\div$ \se  & $1.17\times$ & $0.29\times$ & $0.25$ \\
        \bottomrule
    \end{tabular}%
}
    \end{minipage}
    \caption{E3: could differential measurement reliability explain the ordering? No --- it
    predicts the opposite. \textit{(a)} ICC(2,1) from 75 \sleepedf subjects with two consecutive
    nights; disattenuated (``true'') $r$ divides the observed $r$ by $\sqrt{\mathrm{ICC}}$.
    Correcting for unreliability \emph{widens} the gap against \stage
    ($+0.029$ vs.\ $+0.228$). \textit{(b)} If \stage and its comparator had identical true
    correlations with amyloid, measurement error alone would still make the observed \stage
    correlation $1.86\times$ larger than \tst's. The observed ratio is about eleven times
    smaller than that prediction.}
    \label{tab:reliability}
\end{table}
